% Appendix D -- Certified numerical computation. To be included after \appendix.
% The outputs quoted come from code/appD_certifie.py and from the robertson example of rusty-SUNDIALS, run on 2026-10-26.
% Definitions shared by appendices C, D and E (each carries them, so that it can be included on its own).
\ifdefined\chemin\else
\lstdefinestyle{appC}{language=Lean,
  literate={ℝ}{{$\mathbb{R}$}}1 {ℕ}{{$\mathbb{N}$}}1 {≤}{{$\le$}}1 {→}{{$\to$}}1 {∀}{{$\forall$}}1 {∃}{{$\exists$}}1
  {∧}{{$\wedge$}}1 {∈}{{$\in$}}1 {∑}{{$\sum$}}1 {·}{{$\cdot$}}1 {⟨}{{$\langle$}}1 {⟩}{{$\rangle$}}1 {α}{{$\alpha$}}1
  {é}{{\'e}}1 {è}{{\`e}}1 {ê}{{\^e}}1 {à}{{\`a}}1 {ô}{{\^o}}1 {ç}{{\c{c}}}1 {«}{{\guillemotleft}}1 {»}{{\guillemotright}}1
  {’}{{'}}1}
\DeclareUrlCommand\chemin{\urlstyle{tt}\def\UrlBreaks{\do\/\do\_\do\-\do\.}}
\DeclareUrlCommand\sha{\urlstyle{tt}\def\UrlBreaks{\do0\do1\do2\do3\do4\do5\do6\do7\do8\do9\do a\do b\do c\do d\do e\do f}}
\fi
\chapter{Certified numerical computation}\label{appD:chap}

A number computed in floating point is an approximation whose error is unknown. A \emph{certified} number comes with a
proved enclosure. Between the two, this book distinguishes several degrees (§\ref{ch1:sec:tiers}):
\begin{itemize}
\item a floating-point computation, however precise, is \tierX{};
\item an enclosure in ball arithmetic (Arb) is rigorous \emph{if the code is correct}; since that code is not verified
by Lean, the book classifies it as \tierC{} (``C-ball'' in Chapter~3);
\item a claim about \emph{stated rationals}, decided by Lean's kernel, is \tierAs{} for that arithmetic claim alone, and
\tierA{} when the module has in addition been replayed by \lean{leanchecker} with code~$0$.
\end{itemize}
This appendix explains these three levels, then the solution of stiff differential equations with CVODE, and why a
CVODE result remains Tier~X. A complete example, \chemin{code/appD_certifie.py}, goes through the three levels with
its controls.

%=====================================================================================================================
\section{Floating point}\label{appD:sec:flottant}

A binary64 floating-point number (IEEE~754 standard~\cite{app:ieee754}) is written $\pm m\cdot2^{e}$ with a 53-bit
significand. The result of an operation ($+,-,\times,\div,\sqrt{\ }$) is the rounding of the exact result, and the
standard guarantees
\[\mathrm{fl}(x\circ y)=(x\circ y)(1+\varepsilon),\qquad|\varepsilon|\le u=2^{-53}\approx1.1\times10^{-16},\]
in round-to-nearest (barring overflow). The standard also offers \emph{directed roundings}: towards
$-\infty$ and towards $+\infty$.

\begin{uniremark}\tierL{} The bound above is an instance of the standard model of floating-point arithmetic: for every
real $x$ in the normal range of binary64, round-to-nearest satisfies $\mathrm{fl}(x)=x(1+\varepsilon)$ with
$|\varepsilon|\le u$, i.e.\ $|\mathrm{fl}(x)-x|\le u|x|$; correct rounding of $+,-,\times,\div,\sqrt{\ }$ means
$\mathrm{fl}(x\circ y)$ is $\mathrm{fl}$ applied to the exact $x\circ y$, which gives the displayed bound~\cite{app:ieee754}.
\end{uniremark}

\begin{example}[Ten times one tenth]
$1/10$ is not representable in binary. The floating-point number \texttt{0.1} is exactly
$0.1000000000000000055511\ldots$, slightly more than $1/10$. In Python, \texttt{sum([0.1]*10)} gives
\chemin{0.9999999999999999} and not $1$: the rounding error accumulates, and the sign of the final error is not that
of the initial error.
\end{example}

The elementary functions ($\exp$, $\sin$\ldots) of the system's mathematical library have in general \emph{no}
guarantee of correct rounding: the IEEE~754 standard recommends it without requiring it. A rigorous enclosure therefore
cannot rely on \texttt{math.exp}.

%=====================================================================================================================
\section{Interval arithmetic}\label{appD:sec:intervalles}

\begin{definition}
An unknown real $x$ is represented by an interval $[\underline x,\overline x]$ that contains it. The operations are
defined by
\[[a,b]+[c,d]=[a+c,\,b+d],\qquad[a,b]-[c,d]=[a-d,\,b-c],\]
\[[a,b]\times[c,d]=[\min S,\max S],\quad\text{where }S=\{ac,ad,bc,bd\}.\]
\end{definition}

\begin{theorem}[Inclusion property]\label{appD:thm:inclusion}
If $x\in[a,b]$ and $y\in[c,d]$, then $x+y\in[a,b]+[c,d]$, $x-y\in[a,b]-[c,d]$ and $xy\in[a,b]\times[c,d]$. By
induction, every expression built from these operations, evaluated on intervals containing the arguments, yields an
interval that contains the exact value.
\end{theorem}

\begin{humanproof}
For the sum: $a\le x\le b$ and $c\le y\le d$ give $a+c\le x+y\le b+d$. For the difference, add $-d\le-y\le-c$. For the
product, for fixed $y$ the map $x\mapsto xy$ is affine, so $xy$ lies between $ay$ and $by$; then $y\mapsto ay$ and
$y\mapsto by$ are affine, so $ay$ and $by$ lie between their values at $c$ and $d$. Hence $xy$ lies between the
smallest and the largest of the four products.
\end{humanproof}

\paragraph{Directed rounding.} On a machine, the endpoints are floating-point numbers and every operation rounds. To
preserve the inclusion property, one rounds the lower endpoint \emph{downwards} and the upper endpoint \emph{upwards}:
this is \emph{outward rounding}. Without access to the rounding modes, one can widen each endpoint by one ulp (the gap
between two consecutive floating-point numbers) with \texttt{math.nextafter}: this is wider than necessary, but safe,
since round-to-nearest is off by at most half an ulp.

\paragraph{The dependency problem.} Interval arithmetic does not know that two occurrences of the same variable are
equal: for $x\in[0,1]$, $x-x$ gives $[0,1]-[0,1]=[-1,1]$, and not $\{0\}$. Enclosures grow quickly unless one takes
care; the remedies are to rewrite the expressions, to subdivide the domain (the 206\,468 cells of B2 in Chapter~3) or
to use Taylor expansions.

\begin{numbox}[: step 1 of \chemin{code/appD_certifie.py}]
Ten additions of the interval $[\texttt{nextafter(0.1,-inf)},\,\texttt{0.1}]$, which contains $1/10$, widened by one
ulp at each addition. Actual output (2026-10-26):
\begin{lstlisting}[style=appC,language={}]
"naive_sum": "0.9999999999999999",   "naive_equals_1": false,
"outward": ["0.9999999999999991", "1.0000000000000009"],   "outward_contains_1": true,
"nearest": ["0.9999999999999999", "0.9999999999999999"],   "nearest_contains_1": false
\end{lstlisting}
\emph{Positive control:} the outward-rounded interval contains exactly $1$ (the membership test is done in exact
rationals, \chemin{fractions.Fraction}). \emph{Negative control:} the same computation, with endpoints rounded to
nearest, gives an ``interval'' that excludes the true value.
\end{numbox}

%=====================================================================================================================
\section{Ball arithmetic: Arb and python-flint}\label{appD:sec:boules}

Arb~\cite{ch3-arb} represents a real number by a \emph{ball} $[m\pm r]$: a midpoint $m$ in arbitrary precision (a
dyadic number $M\cdot2^E$) and a radius $r$ in low precision, rounded upwards. It is interval arithmetic in
midpoint--radius form, which is more economical when the precision is high.

\begin{proposition}[Product of two balls]\label{appD:prop:produit}
If $|x-m_1|\le r_1$ and $|y-m_2|\le r_2$, then $|xy-m_1m_2|\le|m_1|r_2+|m_2|r_1+r_1r_2$.
\end{proposition}

\begin{humanproof}
$xy-m_1m_2=m_1(y-m_2)+m_2(x-m_1)+(x-m_1)(y-m_2)$; apply the triangle inequality.
\end{humanproof}
On a machine, Arb computes the rounded midpoint $m_1m_2$, and adds to the radius the error of that rounding, bounded
from above. For functions ($\exp$, $\operatorname{erf}$, integrals\ldots), the library uses expansions with an explicit
bound on the remainder. The result contains the exact value \emph{if Arb's code is correct}: it is tested and
published software, not a formally verified program.

\paragraph{In Python.} The \texttt{python-flint} library gives access to Arb. The B2 study of Chapter~3 used
version~0.6.0, installed in \chemin{.scratchpad/b2_pylib} of the repository for the Python~3.11 of the \texttt{.venv}
environment; the system \texttt{python3} (3.10) carries version~0.9.0 instead (Appendix~\ref{appE:chap}). The example
in this appendix uses version~0.6.0.

\begin{numbox}[: step 2 of \chemin{code/appD_certifie.py}]
Enclosure of $I=\int_0^1e^{-x^2}dx$ by \texttt{acb.integral} at 96~bits, checked against the independent ball
$\frac{\sqrt\pi}2\operatorname{erf}(1)$. Actual output:
\begin{lstlisting}[style=appC,language={}]
"python_flint": "0.6.0",  "prec_bits": 96,
"ball":     "[0.746824132812427025399467436 +/- 2.30e-28]",
"erf_ball": "[0.7468241328124270253994674361 +/- 5.34e-29]",   "erf_overlaps": true,
"midpoint_float_n10": "0.7471308777479975",   "midpoint_inside_ball": false
\end{lstlisting}
\emph{Positive control:} the two balls, computed along two routes (rigorous quadrature, error function), overlap.
\emph{Negative control:} the $10$-point midpoint rule, in floating point, gives $0.74713$: a plausible value, wrong in
the fourth digit, and outside the ball. A floating-point number is not an enclosure.
\end{numbox}

%=====================================================================================================================
\section{Rational certificates checked by Lean: LeanFlow}\label{appD:sec:leanflow}

\subsection{The principle}
Lean's kernel can decide equalities and inequalities between explicit \emph{integers}: this is a finite computation.
LeanFlow (\texttt{leanflow\_mcp.lean}, function \texttt{certify}) turns a numerical claim into such a statement, in four
steps:
\begin{enumerate}
\item \textbf{The claim.} Named rational enclosures, \texttt{x = [lo, hi]}, and linear inequalities between their
endpoints. Numbers are written \texttt{p/q} or in decimal; a binary floating-point number is rejected (it is Tier X).
\item \textbf{The template.} Each inequality becomes a theorem over $\Z$, with denominators cleared symbolically, of the
form \texttt{theorem cert\_i : (0 : Int) < D $\wedge$ N * RD $\le$ RN * D := by decide +kernel}. Python only assembles
text from integer literals; all multiplications and comparisons are done by Lean.
\item \textbf{The kernel decides.} The tactic \lean{decide +kernel} (LeanFlow name: \texttt{decide\_kernel}) has the
proposition evaluated by the kernel itself, with no additional axiom.
\item \textbf{The audit.} LeanFlow appends to the source one \texttt{\#print axioms} per theorem and a random sentinel,
then requires exactly one report per theorem, within the whitelist, and the sentinel at the end. An exact oracle in
Python (\texttt{Fraction}) must return the same verdict.
\end{enumerate}

\paragraph{Fail-closed.} ``Closed when in doubt'': every abnormal situation yields \texttt{error}, never \texttt{ok}.
LeanFlow's documentation (\chemin{docs/LEAN_CERTIFY.md}) lists the cases, each with a test: missing Lean toolchain,
imported module without \texttt{.olean}, nonzero return code, no axiom report, \lean{sorry}, user axiom,
\lean{native_decide} (which introduces an axiom specific to the use), forged axiom report, \texttt{\#exit}, invalid
claim, disagreement with the oracle. A \emph{false} claim yields \texttt{rejected} with the code
\texttt{false-statement}.

\paragraph{What Tier A covers.} The status \texttt{ok} concerns the arithmetic statement \emph{about the stated
numbers}, and nothing else: neither the provenance of these numbers, nor the fact that they really enclose the
integral. This is why, in Chapter~3, the B2 chain is \tierC{} overall: the arithmetic is verified, the ball is not.
LeanFlow audits only \texttt{\#print axioms}; the book adds the \lean{leanchecker} replay when it claims \tierA{}
(Appendix~\ref{appC:chap}, §\ref{appC:sec:leanchecker}, and case U5 of Chapter~5).

\subsection{The complete example}
The ball of step~2 is rounded outwards to two decimals with 30~digits, in exact rationals:
$\mathtt{lo}=\lfloor10^{30}(m-r)\rfloor/10^{30}$ and $\mathtt{hi}=\lceil10^{30}(m+r)\rceil/10^{30}$. LeanFlow then
receives two assertions:
\begin{lstlisting}[style=appC,language={}]
{"name": "gaussI", "tactic": "decide_kernel",
 "enclosures": {"I": ["186706033203106756349866858969/250000000000000000000000000000",
                      "74682413281242702539946743623/100000000000000000000000000000"]},
 "assertions": [{"label": "I.lo > 0.7468", "endpoint": "I.lo", "op": ">", "rhs": "7468/10000"},
                {"label": "I.hi < 0.7469", "endpoint": "I.hi", "op": "<", "rhs": "7469/10000"}]}
\end{lstlisting}
and, as a negative control, the same ball with the false assertion \texttt{I.hi < 7468/10000}.
\begin{numbox}[: step 3 of \chemin{code/appD_certifie.py}]
Actual output (Lean 4.34.1, commit \sha{5045d0056413266e57c625dcd7c365b10e377c52}, core only):
\begin{lstlisting}[style=appC,language={}]
"pos_I_between_0.7468_and_0.7469": {"status": "ok", "code": "verified", "tier": "A", "failed": [],
                                    "seconds": 0.64}
"neg_I_below_0.7468":              {"status": "rejected", "code": "false-statement", "tier": null,
                                    "failed": ["cert_0"], "seconds": 0.49}
"all_controls_pass": true
\end{lstlisting}
The script's seven controls (three for step~1, two for each of steps~2 and~3) passed; the whole script took 1.7~s.
\end{numbox}
\noindent Tier summary: the statement ``$0.7468<\mathtt{lo}$ and $\mathtt{hi}<0.7469$'' is \tierAs{}: LeanFlow
returns \texttt{"tier": "A"}, but we did not replay these certificates with \lean{leanchecker}, and the book's
convention reserves \tierA{} for modules replayed with code~$0$; the claim ``$I\in[\mathtt{lo},\mathtt{hi}]$'' is
\tierC{} (it rests on Arb); the conclusion ``$0.7468<I<0.7469$'' inherits the weaker of the two, \tierC.

%=====================================================================================================================
\section{Stiff differential equations and CVODE}\label{appD:sec:cvode}

\subsection{Stiffness}
Consider $y'=-\lambda y$ with $\lambda>0$ large. The solution $e^{-\lambda t}$ is trivial, but the explicit Euler
method $y_{n+1}=y_n+hf(y_n)=(1-h\lambda)y_n$ stays bounded only if $|1-h\lambda|\le1$, that is, $h\le2/\lambda$. A system
is \emph{stiff} when it mixes very different time scales: the stability of an explicit method then imposes a step
dictated by the fastest one, much smaller than what accuracy on the slowest one requires. The \emph{implicit} Euler
method, $y_{n+1}=y_n+hf(y_{n+1})$, gives $y_{n+1}=y_n/(1+h\lambda)$: it is stable for every $h>0$.

\subsection{BDF formulas}
CVODE~\cite{app:sundials} integrates stiff systems with the \emph{backward differentiation formulas} (BDF) of order
$q=1,\dots,5$. One approximates $y'(t_{n+1})$ by the derivative of the polynomial interpolating
$y_{n+1},y_n,\dots,y_{n+1-q}$:
\[\text{BDF1: }y_{n+1}-y_n=hf(y_{n+1}),\qquad\text{BDF2: }y_{n+1}-\tfrac43y_n+\tfrac13y_{n-1}=\tfrac23hf(y_{n+1}).\]
Each step requires solving a \emph{nonlinear} equation in $y_{n+1}$, of the form $G(y)=y-\gamma hf(y)-a_n=0$, where
$a_n$ collects the past values and $\gamma$ depends on the order.

\subsection{Newton and the Jacobian}
CVODE solves $G(y)=0$ by a modified Newton method:
$\bigl(I-\gamma hJ\bigr)\,\delta^{(k)}=-G(y^{(k)})$, $y^{(k+1)}=y^{(k)}+\delta^{(k)}$, where $J=\partial f/\partial y$ is the
Jacobian. The matrix $I-\gamma hJ$ is recomputed and factorised only from time to time. The Jacobian is either supplied
analytically by the user or approximated by finite differences, column by column (one extra evaluation of $f$ per
column). The book's \texttt{qf-*} crates supply analytic Jacobians: for a gradient flow $\dot x=-\nabla E(x)$
(relaxations in cases U1 and U2), $J=-\mathrm{Hess}\,E$; for the Kolmogorov equation $\dot\pi=\pi(P-I)$ (case U5), the
system is linear and $J$ is exact.

\subsection{Error control, and what it does not guarantee}
At each step, CVODE estimates the \emph{local} error $e$ (the error made in one step, starting from values assumed
exact) and accepts the step if
\[\|e\|_{\mathrm{WRMS}}=\Bigl(\frac1n\sum_{i=1}^n\Bigl(\frac{e_i}{\mathrm{rtol}\,|y_i|+\mathrm{atol}_i}\Bigr)^2\Bigr)^{1/2}\le1;\]
otherwise it reduces the step. It also adjusts the order $q$. This control is \emph{heuristic}: the global error, after
thousands of steps, is not bounded by rtol; the local error estimate rests on asymptotic expansions; and everything is
done in floating point. This is why \textbf{a CVODE result is Tier~X}. The book never presents it as a proof: it
compares it with an exact value (U5), with a conservation law (U4), or with a theorem (U1, U2), with controls.

\paragraph{The \texttt{qf-popgen} certificate is not a proof.} The ``absorption certificate'' of case U5 (stochastic
kernel, conserved mass, no negative component, interior mass $<10^{-9}$, final value within $10^{-8}$ of the target) is
a set of \emph{checks} that reject dubious runs. It makes a run more reliable; it does not prove it. The only Tier~A
part of case U5 is the LeanFlow statement about the number \chemin{0.5714285714284759} and about the vector
$(0,\frac47,\frac67,1)$.

\subsection{A stiff example: Robertson kinetics}
The Robertson problem is the classic stiffness test, with rate constants from $0.04$ to $3\times10^{7}$:
\[y_1'=-0.04y_1+10^4y_2y_3,\quad y_2'=0.04y_1-10^4y_2y_3-3\times10^7y_2^2,\quad y_3'=3\times10^7y_2^2,\]
with $y(0)=(1,0,0)$ and $t\in[0,4\times10^{10}]$. The sum $y_1+y_2+y_3$ is conserved exactly by the system, which
provides a check. rusty-SUNDIALS ships it as the example \chemin{examples/robertson.rs} (BDF up to order $5$,
$\mathrm{rtol}=10^{-4}$, $\mathrm{atol}=10^{-8}$, analytic Jacobian). Command, from the working copy
\chemin{rusty-SUNDIALS-wt-tb2ch5}:
\begin{lstlisting}[style=appC,language={}]
cargo run --manifest-path ~/rusty-SUNDIALS-wt-tb2ch5/Cargo.toml --target-dir ~/rusty-target-ar \
    --release -j 2 --offline -p examples --example robertson
\end{lstlisting}
\begin{numbox}[: actual output (2026-10-26, 43~s including compilation)]
\begin{lstlisting}[style=appC,language={}]
           t             y1             y2             y3
----------------------------------------------------------
   4.0000e-1    9.851721e-1    3.386396e-5    1.479401e-2
    4.0000e0    9.055204e-1    2.240497e-5    9.445716e-2
    4.0000e2    4.505200e-1    3.222938e-6    5.494767e-1
    4.0000e6    5.169270e-4    2.068765e-9    9.994831e-1
   4.0000e10    5.155066e-8   2.062027e-13    9.999999e-1

Solver statistics:
  Steps:      633  (C ref: 1070, ratio: 0.6x)
  RHS evals:  1762  (C ref: 1537, ratio: 1.1x)
  Final order: 5
  Conservation (y1+y2+y3): 1.000000000000000e0 (should be 1.0)
  Conservation error: 4.44e-16  (threshold 1e-12)
\end{lstlisting}
(Intermediate lines omitted; the program's last line, \texttt{Conservation: PASS} followed by a check mark, is cut
here because the check mark is not printable.)
\end{numbox}
\noindent Reading: in $633$ steps, CVODE traverses eleven orders of magnitude in time; the step size grows accordingly.
An explicit method would be limited in stability, over the whole interval, by the fast eigenvalues of the Jacobian. The
program compares its statistics with the C version of SUNDIALS (\texttt{cvRoberts\_dns}, reference values written in
the example's code, not rerun here; with an absolute tolerance different from that of the Rust example, so the comparison of step counts is only indicative). Conservation to within $4.4\times10^{-16}$ is a \emph{necessary} check; it proves
nothing about $y_2$, which could be wrong without violating conservation. This is a \tierX{} computation.

%=====================================================================================================================
\section{Exercises}\label{appD:sec:exos}

\begin{exercise}
Prove that, for $x\in[a,b]$ with $0<a$, the interval $[1/b,1/a]$ contains $1/x$. What should be done if $0\in[a,b]$?
\end{exercise}

\begin{exercise}
Evaluate $f(x)=x^2-2x$ on $[0,2]$ in interval arithmetic, then the form $f(x)=(x-1)^2-1$. Compare with the exact image
$[-1,0]$ and explain the discrepancy by the dependency problem.
\end{exercise}

\begin{exercise}
For the BDF2 method applied to $y'=-\lambda y$, write the linear recurrence satisfied by $y_n$ and show that its
characteristic roots have modulus $<1$ for every $h\lambda>0$ (A-stability of BDF2).
\end{exercise}

\begin{exercise}[LeanFlow, negative control]
In \chemin{code/appD_certifie.py}, change the precision \texttt{ctx.prec} to $20$~bits. Does the ball still contain the
value? Do the two LeanFlow certificates still pass? Why does a wider enclosure make the assertion
\texttt{I.lo > 0.7468} harder to certify?
\end{exercise}
